\documentclass[11pt,a4paper]{article}
\usepackage[margin=2.6cm]{geometry}
\usepackage{amsmath,amssymb,amsthm,mathtools,microtype,booktabs,array,tabularx}
\usepackage[shortlabels]{enumitem}
\usepackage[colorlinks=true,linkcolor=black,urlcolor=blue]{hyperref}
\hypersetup{pdftitle={The two-linear-gate case: a human proof and exhaustive checks},pdfauthor={AI-written draft, pending human review}}
\newtheorem{theorem}{Theorem}
\newtheorem{lemma}[theorem]{Lemma}
\newtheorem{proposition}[theorem]{Proposition}
\newcommand{\F}{\mathbb F_2}
\newcommand{\Aff}{\mathrm{Aff}}
\newcommand{\Span}{\operatorname{span}}
\newcommand{\rk}{\operatorname{rank}}
\newcommand{\supp}{\operatorname{supp}}
\newcommand{\rad}{\operatorname{rad}}
\newcommand{\status}[1]{\hfill\textsf{[#1]}}
\setlength{\parskip}{3pt}
\title{The two-linear-gate case of the low-half multiplier:\\a human proof and exhaustive checks}
\author{AI-written draft, pending human review}
\date{1 November 2026}
\begin{document}
\maketitle
\begin{abstract}
We give a human proof that five AND gates with exactly two affine-input gates cannot compute the three low bits of a three-bit product. The new step is a product lemma modulo the Boolean ideal $(a_0b_0)$: a quadratic product in that quotient introduces at most one rank-two direction. Combined with multiplier rigidity and an exterior-support obstruction, it excludes the remaining higher-degree branch. Inside the four-variable core we use the given coset lemmas from the three-linear-gate proof. Every new lemma has an exhaustive check of its full application domain, and an independent complete circuit search reproduces the 2,064 previously difficult pairs. The draft awaits human review.
\end{abstract}

\section{Setting and imports}
The inputs, in order, are $a_0,a_1,a_2,b_0,b_1,b_2$. Boolean functions are their ANFs, with $x_i^2=x_i$. Put
\[
\Aff=\Span\{1,a_0,a_1,a_2,b_0,b_1,b_2\},\qquad
V(Q)=\{f:\deg f\le2,\ f_2\in Q\}.
\]
The targets are
\[
z_0=a_0b_0,\quad z_1=a_0b_1+a_1b_0,\quad
z_2=a_0b_2+a_1b_1+a_2b_0+a_0a_1b_0b_1.
\]
Write $P_0=a_0\wedge b_0$, $Z_1=(z_1)_2$, $\Omega=a_0\wedge a_1\wedge b_0\wedge b_1$, $U=\langle a_0,a_1,b_0,b_1\rangle$, and $E=\{e_{a_2},e_{b_2},e_{a_2}+e_{b_2}\}$. A gate is linear when both factors, as functions, are affine.

The supplied proofs are treated as GIVEN and are not re-proved: [L4, Lemma 1] (normal form), [L4, Lemma 2] (quadratic spans), [L4, Lemma 3] (support), [L4, Lemma 8] (derivatives), [L4, Lemma 11] (types), [L4, Theorem 12]; and [L3, Lemma P (Lemma 6)], [L3, Radical lemma (Lemma 4)], [L3, Weights lemma (Lemma 5)], and [L3, Theorem 3]. Source labels are also given because the supplied files have no compiled cross-reference files.

Thus the normalized prefix is $g_0=a_0b_0$, one more affine-input gate with rank-two quadratic part $P$, and then $g_2,g_3,g_4$. Put $Q=\Span\{P_0,P\}$. Dependent or zero quadratic parts are allowed by taking $\dim Q\le2$; a redundant initial gate can instead be replaced by an independent affine-input gate, enlarging the available space.

\section{The theorem and the plan}
\begin{theorem}\label{thm:main}\status{PROVED}
No normalized five-AND circuit with exactly two linear gates computes $(z_0,z_1,z_2)$.
\end{theorem}
We first dispose of $\deg g_2\le2$, then exclude every space obtained by adjoining one product to $V(Q+\langle Z_1\rangle)$. This also excludes early placements of $z_2$ or $z_1+z_2$.

\begin{lemma}[Replacing a quadratic nonlinear gate]\label{lem:replace}\status{PROVED}
In the normalized setting, let $Q$ contain $P_0$ and be generated by at most $m$ independent rank-two forms, $m\in\{1,2\}$, and let $f,h\in V(Q)$, with at least one nonaffine. If $\deg(fh)\le2$, there is a rank-two-generated $(m+1)$-dimensional space $Q^+$ containing $Q$ such that $fh\in V(Q^+)$. A circuit prefix of $m$ linear gates followed by this product can be replaced by $m+1$ linear gates, preserving availability of every original prefix function and every subsequent factor function.
\end{lemma}
\begin{proof}
By [L3, Lemma P (Lemma 6)], $(fh)_2=\pi+\rho$, with $\pi\in\Span\{f_2,h_2\}\subseteq Q$ and either $\rho=0$ or $\rk\rho=2$. If $\rho=l\wedge r\ne0$, the affine-input gate $lr$ supplies precisely this new quadratic direction; its affine part is free. Therefore $fh\in V(Q+\langle\rho\rangle)$ by [L4, Lemma 2]. Extend a rank-two generating basis of this space to dimension $m+1$ using quadratic monomials. This is possible because all quadratic monomials have rank two and span the 15-dimensional quadratic space. In particular, if $\rho=0$ or $\rho\in Q$, choose any new independent rank-two direction. If the original prefix has dependent gates, replace redundant gates during this extension.

Use the resulting $m+1$ affine-input products as the new prefix. Its available space contains $V(Q)+\langle fh\rangle$. Keep every later factor as the same Boolean function as before, expressing it in the enlarged space. Induction through the tail preserves all its signals. A factor that was nonaffine remains nonaffine, since its function has not changed.
\end{proof}
For $m=2$ the replacement has three initial linear gates and the original two nonlinear tail gates. It is already in the form covered by [L3, Theorem 3]; independent quadratic directions have been supplied even in the degenerate cases. Consequently the branch $\deg g_2\le2$ is excluded by a human proof.

\begin{lemma}[The 32 symmetries]\label{lem:sym}\status{PROVED}
Interchanging the two operands and independently making the substitutions
\[
a_2\mapsto a_2+s,\quad b_2\mapsto b_2+t,
\qquad s,t\in\langle a_0,b_0\rangle,
\]
generate a group of order 32. They fix $z_0,z_1$ and preserve $z_2$ modulo $\Aff+\langle z_0\rangle$. They preserve semantic linearity, circuit size, and all prefix-space feasibility conditions used here.
\end{lemma}
\begin{proof}
There are 16 distinct commuting shifts. Operand interchange conjugates this shift subgroup to itself and is outside it, giving 32 distinct invertible linear substitutions. A shift changes $z_2$ by $b_0s+a_0t$, which belongs to $\Span\{a_0,b_0,a_0b_0\}$ because $a_0^2=a_0$, $b_0^2=b_0$. Operand interchange fixes the targets. Substitution commutes with products and sums, fixes $P_0$ as a function, and bijects affine functions; all asserted circuit properties follow. For the weaker derivative-and-$\Omega$ diagnostic we make no claim of invariance when only one external direction is tested; that diagnostic is run on the stated representatives and is not used to justify the exhaustive circuit reduction.
\end{proof}

\begin{lemma}[Top degree]\label{lem:degree}\status{PROVED}
For $f,h\in V(Q)$ one has $\deg(fh)\le4$. If $\dim Q=2$ and $\deg(fh)=4$, then $\Span\{f_2,h_2\}=Q$ and $(fh)_4=P_0\wedge P$. If $Q=\langle P_0\rangle$, then $\deg(fh)\le3$.
\end{lemma}
\begin{proof}
The quartic part is $f_2\wedge h_2$. The wedge square of every two-form is zero in characteristic two. Thus a nonzero wedge in a two-dimensional space requires a basis, and all bases over $\F$ have determinant one. In a one-dimensional space the wedge is always zero.
\end{proof}

\section[A structural restriction on feedback placing z1]{A structural restriction on feedback placing $z_1$}
Define the four-element set
\[
\mathcal K=\{1+P_0,\ 1+a_0+P_0,\ 1+b_0+P_0,\ a_0+b_0+P_0\}.
\]
Here and below a two-form in a Boolean expression denotes its quadratic ANF representative; the four functions in $\mathcal K$ are exactly the complements of the indicators of the four fibers of $(a_0,b_0)$.

\begin{lemma}[Four multipliers]\label{lem:four}\status{PROVED}
Suppose either $Q=\langle P_0\rangle$, or $Q=\langle P_0,P\rangle$ with $P$ rank two and $P\notin\Lambda^2(U)$. For $k\in V(Q)$, a function $R\in z_1+V(Q)$ satisfying $kR=R$ exists if and only if $k\in\{1\}\cup\mathcal K$.
\end{lemma}
\begin{proof}
Every quadratic part $R_2=Z_1+q$, $q\in Q$, has rank at least four. Otherwise, in the two-dimensional case, a rank-two $R_2$ would split $Z_1$ or $Z_1+P_0$ into $P$ plus another rank-two form, contradicting [L4, Lemma 3]. The one-dimensional case is immediate.

If $k$ is affine and nonconstant, [L3, Lemma P, part P1] applied to $kR=R$ makes $R_2$ rank at most two. Thus an affine $k$ must be 1. If $k_2$ has rank four, the Radical lemma gives $\supp k_2=\supp R_2$, of dimension four. If $R_2$ has no $P$ coefficient, this support is $U$, impossible for $k_2=P+\epsilon P_0$. If it has a $P$ coefficient, $R_2+k_2=Z_1+\epsilon P_0$ has support $U$, again forcing the common four-dimensional support to be $U$, with the same contradiction.

It remains that $k_2$ has rank two. If $D_ek=1$ for some $e\in\rad k_2$, differentiate $kR=R$. Since $D_eR$ is affine, the derivative identity gives $R=kD_eR$. Part P1 of the product lemma again gives rank at most two (the constant derivative cases are immediate). Therefore $k$ is constant on cosets of $\rad k_2$ and depends only on two independent linear forms $l,r$. A rank-two quadratic function on two bits has weight one or three; on six bits its weight is 16 or 48.

A quadratic function of rank at least four has weight at least 24: [L3, Lemma 5] gives 24 or 40 in rank four and 28 or 36 in rank six; an extra linear term in a radical direction gives weight 32 by pairing. Since $\supp R\subseteq\supp k$, $k$ cannot have weight 16. Thus $1+k$ is the indicator of an affine codimension-two flat, on which $R$ vanishes. In coordinates defining this flat, the quadratic part of $R$ is $l\wedge A+r\wedge B$, of rank at most four. Hence $R_2$ has rank exactly four.

The Radical lemma now gives $\supp k_2\subseteq\supp R_2$. If $k_2$ has a $P$ coefficient, either $R_2$ has none and its support is $U$, or $R_2+k_2=Z_1+\epsilon P_0$ has support $U$. Both alternatives force $\supp R_2=U$ and $k_2\in\Lambda^2(U)$, a contradiction. Consequently $k_2=P_0$. Its weight-48 functions with no radical dependence are exactly $\mathcal K$.

For sufficiency, let $(s,t)$ be the unique zero fiber of $k\in\mathcal K$. Then $R=(a_0+s)b_1+(b_0+t)a_1$ belongs to $z_1+\Aff$, vanishes on that fiber, and satisfies $kR=R$. The case $k=1$ is immediate.
\end{proof}

If a gate using a previous signal $g$ places $z_1$ from $V(Q)+\langle g\rangle$, plane operations and a constant change in its other factor put the placement in the form
\[
k(g+F)\in z_1+V(Q),\qquad k,F\in V(Q).
\]
Indeed replacing two $g$-containing factors by their sum removes $g$ from one factor, and replacing $k$ by $k+1$ absorbs a $g$ term in the output correction. Multiplying the displayed identity by $k$ supplies the hypothesis of Lemma~\ref{lem:four}. If $g$ has degree at least three, $k=1$ is impossible, so only $\mathcal K$ remains.

\section[Rigidity of a quadratic multiplier placing z2]{Rigidity of a quadratic multiplier placing $z_2$}
\begin{lemma}[Last-factor rigidity]\label{lem:lastfactor}\status{PROVED}
Let $P$ have rank two and lie outside $\Lambda^2(U)$, and put $S=\langle P_0,P,Z_1\rangle$. If $k\in V(S)$, $R\in z_2+V(S)$ and $kR=R$, then $k\in\{1,1+P_0\}$.
\end{lemma}
\begin{proof}
Write $k_2=dP+u$ with $u\in\langle P_0,Z_1\rangle$ and $d\in\F$. The quartic part of $R$ is $\Omega$. If $e\in E$ kills $k_2$ and $D_ek=1$, differentiation gives $R=kD_eR$. But $D_eR$ is affine, so the right side has degree at most three, a contradiction. We will use this observation repeatedly.

First suppose $P$ has type $v$ in [L4, Lemma 11], and choose its killing direction $e$. Put $r=\delta_e\in\langle a_0,b_0\rangle\setminus\{0\}$. Since $e$ kills $S$, the observation gives $D_ek=0$. Differentiating $kR=R$ yields $kA=A$, where $A=D_eR=r+c$ for a constant $c$. Therefore
\[
1+k=(1+A)B
\]
for an affine $B$: in coordinates with $A$ as one variable, comparison on $A=1$ proves this factorization and the degree bound on $B$. In particular $k_2$ has rank at most two and, if nonzero, has factor $r$.

If $d=0$, the only nonzero rank-two form in $\langle P_0,Z_1\rangle$ is $P_0$. The factorization then makes $k$ a function of $a_0,b_0$ alone, including the affine and constant cases. Suppose it vanished on a fiber $(s,t)\ne(1,1)$. The quadratic part of $R$ on that fiber, in the four remaining variables, would be
\[
a_1\wedge b_1+\epsilon\overline P.
\]
This cannot be zero: $\overline P=a_1\wedge b_1$ would imply $P\in\Lambda^2(U)$ by projecting the two-dimensional support of the decomposable form $P$. Thus $k$ is 1 on all three such fibers, giving precisely $k=1$ or $k=1+P_0$.

Suppose $d=1$. A $Z_1$ coefficient in $k_2$ is impossible by [L4, Lemma 3], applied to the two rank-two forms $P,k_2$ whose sum would be $Z_1$ or $Z_1+P_0$. Hence $k_2=P+\tau P_0$. If $\tau=1$, decomposability of the sum forces $P$ and $P_0$ to share a factor. In either case the $U$-line in the support of $P$ must be $\langle r\rangle$, since $r\in\supp k_2$. Consequently $P=r(p+x)$ with $x\in U$ and a nonzero external linear form $p$ killed by $e$.

Choose another external direction $w$ and its dual coordinate $q$ so that $p(w)=1$, $q(e)=1$, $q(w)=0$. Put $s=\delta_w$, independent of $r$. The cross part of $q_2$ is $rq+sp$. In $1+k=(1+r+c)B$, the affine function $B$ has $p$ coefficient one and $q$ coefficient zero. On the codimension-two flat
\[
r=c,\qquad B=1,
\]
we have $k=0$, so $R=0$; the second equation expresses $p$ affinely in the remaining variables $s,a_1,b_1$.

On $r=c$, $a_0b_0$ is either zero or a nonconstant affine function of $s$. In the latter case $\Omega$ retains a cubic term $sa_1b_1$, while all other terms of $R$ remain quadratic after the affine substitution for $p$. This is impossible. In the former case the carry vanishes. The remaining quadratic part on the flat is then $a_1\wedge b_1+s\wedge L$ for some linear $L$: $rq$ is affine, $sp$ has factor $s$, $P=r(p+x)$ becomes affine, and $Z_1$ restricts to a quadratic with factor $s$. Its $a_1b_1$ coefficient is one, again impossible. This excludes $d=1$ in the type case.

It remains that $P$ is full. By [L4, Lemma 11], shift external coordinates by linear forms in $U$ so that $P=pq$. These auxiliary coordinate changes are used only to analyze the identity; they need not preserve the target. They fix $U$, $\Omega$, $P_0$, and $Z_1$, and leave the external coefficients of $q_2$ equal to $b_0,a_0$. If $d=1$, further constant shifts remove the external affine terms of $k$, giving
\[
k=pq+K(u),\quad R=\epsilon pq+B(u)p+A(u)q+T(u),
\]
where $\deg K\le2$, $A,B$ are affine, and $T_4=\Omega$. The $pq$ coefficient in $kR=R$ reads $A+B+T+\epsilon K=0$, contradicting $T_4=\Omega$.

If $d=0$, the initial observation applied to both external unit directions makes $k$ independent of $p,q$. If $\epsilon=1$, comparison of $pq$ coefficients gives $k=1$. If $\epsilon=0$, the $p,q$ coefficients give $(1+k)A=(1+k)B=0$, with $A=a_0+c_0$, $B=b_0+c_1$. A quadratic function supported on the flat $A=B=0$ is a constant multiple of $(1+A)(1+B)$: comparison in coordinates $A,B,a_1,b_1$ proves this directly. Thus $k$ depends only on $a_0,b_0$. The same fiber argument as above leaves only $1$ and $1+P_0$.
\end{proof}

\begin{lemma}[Two feedback placements cannot coexist]\label{lem:two}\status{PROVED}
Keep $Q=\langle P_0,P\rangle$, $S=Q+\langle Z_1\rangle$, with $P$ rank two outside $\Lambda^2(U)$. For any Boolean function $g$, there are no $F,v\in V(Q)$, $H,w\in V(S)$ and $k\in\mathcal K$ such that
\[
k(g+F)=z_1+v,\qquad (1+P_0)(g+H)=z_2+w.
\]
No degree hypothesis on $g$ is required.
\end{lemma}
\begin{proof}
Choose a fiber $(a_0,b_0)=(s,t)\ne(1,1)$ on which $k=1$. Such a fiber exists because $k$ vanishes on only one of the four fibers. Both displayed multipliers are one on the chosen fiber. The first equation makes $g$ belong there to $\Aff_4+\langle\overline P\rangle$, since $z_1$ restricts to an affine function. The second makes it belong to $a_1b_1+\Aff_4+\langle\overline P\rangle$, since the carry vanishes and the external terms of $z_2$ become affine. Therefore $\overline P=a_1\wedge b_1$.

Write $P=l\wedge r$. Its projections to $\langle a_1,a_2,b_1,b_2\rangle$ span $\langle a_1,b_1\rangle$, because their wedge is $a_1\wedge b_1\ne0$. Thus $l,r\in U$, contradicting $P\notin\Lambda^2(U)$.
\end{proof}
\emph{Remark.} Lemma~\ref{lem:two} and Proposition~\ref{prop:reorder} are not used in the proof of Theorem~\ref{thm:main} below, which goes through Proposition~\ref{prop:onehigh}; Lemma~\ref{lem:four} is used only in the companion one-linear-gate draft. They are kept as independent confirmations. Together Lemmas~\ref{lem:four}--\ref{lem:two} exclude the branch where $z_1$ is placed by a gate using $g_2$, and the last gate also uses $g_2$, provided $z_1$ has already been placed. Indeed the space after placement is $V(S)+\langle g_2\rangle$; normalization gives the two displayed equations. The constant multiplier 1 at the final placement would put $g_2$ in $z_2+V(S)$, excluded by [L3, Lemma C (Lemma 9)] after extending $Q$ to three rank-two directions. At the first placement it would make $g_2$ quadratic, also impossible.

\begin{proposition}[Rank-two reordering]\label{prop:reorder}\status{PROVED}
If $\deg g_2\ge3$, $z_1\in V_4$, and $Z_1+q$ has rank two for some $q\in Q$, then the five-gate circuit is excluded by [L3, Theorem 3].
\end{proposition}
\begin{proof}
Here $z_1\notin V(Q)+\langle g_2\rangle$, so $V_4=V(Q)+\langle g_2,z_1\rangle$. Insert a third linear gate with quadratic part $Z_1+q$, delete the old $g_3$, and move the old $g_2$ after the three linear gates. Its original factors remain in $V(Q)$ and are unchanged functions. The new space before $g_4$ is exactly the old $V_4$, so the original last gate remains valid and nonlinear. This is a three-linear-gate circuit in the imported normal form, a contradiction.
\end{proof}

\section{Quadratic extraction modulo the low-bit ideal}
Let $I=(a_0b_0)$ be the ideal in the Boolean function algebra generated by the function $a_0b_0$. In the original ANF coordinates, $I$ consists exactly of the functions whose monomials all contain $a_0b_0$. A quadratic function in $I$ is therefore a scalar multiple of $a_0b_0$.

\begin{lemma}[Quadratic extraction modulo $I$]\label{lem:ideal}\status{PROVED}
Let $Q=\langle P_0,P\rangle$ with $P$ rank two, allowing $Q=\langle P_0\rangle$. If $f,h\in V(Q)$ and a quadratic function $R$ satisfies $fh\equiv R\pmod I$, then $R_2=\pi+\rho$ with $\pi\in Q$ and $\rk\rho\le2$.
\end{lemma}
\begin{proof}
If $\dim Q=1$, both factors are affine modulo $I$, so their product has rank-at-most-two quadratic part. Any other quadratic representative differs by $P_0$.

Now $\dim Q=2$. Choose homogeneous linear $l,r$ with $P=l\wedge r$, and use the actual product $p=lr$ as representative. Its affine difference from the pure quadratic ANF of $P$ is free. Modulo $I$, the factors belong to $\Aff+\langle p\rangle$. If both are affine the result is immediate. Otherwise plane operations put them in the form $p+A$ and $B$, with $A$ affine and $B$ homogeneous linear, changing the product only by an available quadratic function. A constant in $B$ can also be removed by such a change. The case $B=0$ is immediate.

The cubic part of $pB+AB$ is $P\wedge B$. Quadraticity modulo $I$ forces it to equal $P_0\wedge L$ for some linear $L$. If it is zero, then $B\in\langle l,r\rangle$, so $pB$ is a scalar multiple of $p$, using $l^2=l$ and $r^2=r$. The new quadratic direction is $A_1\wedge B$.

If this cubic form is nonzero, the decomposable three-forms $P\wedge B=P_0\wedge L$ have the same three-dimensional support $W$. The distinct two-dimensional supports of $P$ and $P_0$ are contained in $W$ and meet in a line $\langle t\rangle$. Choose $L\in W$. Then
\[
C=pB+(a_0b_0)L
\]
has degree at most two and depends only on coordinates in $W$, so $C_2\in\Lambda^2(W)$. Also $Q=t\wedge W$, of dimension two in the three-dimensional space $\Lambda^2(W)$.

Since $B\notin\supp P$, it is not in $\langle t\rangle$. Complete $t,B$ to a basis $t,B,D$ of $W$; then $B\wedge D\notin t\wedge W$. Thus $B\wedge W$ maps onto $\Lambda^2(W)/Q$. Choose $D'\in W$ with $C_2+B\wedge D'\in Q$. Modulo $I$ the product is $AB+C$, whose quadratic part modulo $Q$ is
\[
B\wedge(A_1+D'),
\]
of rank at most two. Restoring the plane operations changes quadratic parts only in $Q$, and changing quadratic representatives modulo $I$ changes them only by $P_0$. This proves the assertion for the original factors and $R$.
\end{proof}

\begin{lemma}[Rank and exterior obstructions]\label{lem:rankext}\status{PROVED}
For every $\epsilon,\alpha\in\mathbb F_2$, the form $q_2+\epsilon Z_1+\alpha P_0$ has rank six. Adding any rank-two form leaves rank at least four. If $P$ has rank two and lies outside $\Lambda^2(U)$, then $P\wedge\theta\ne\Omega$ for every two-form $\theta$.
\end{lemma}
\begin{proof}
In the radical equations the $b_2,a_2$ rows first force the $a_0,b_0$ coordinates to vanish. The $a_1,b_1$ rows then force the $b_1,a_1$ coordinates to vanish, and the $a_0,b_0$ rows finally force $b_2,a_2$ to vanish. This is independent of $\epsilon,\alpha$. The rank inequality $\rk(A+B)\ge\rk A-\rk B$ gives the second assertion.

Write $P=l\wedge r$. If $P\wedge\theta=\Omega$, then $l\wedge\Omega=r\wedge\Omega=0$. The linear forms $s$ with $s\wedge\Omega=0$ are exactly $U$, as expansion in a basis extending its four factors shows. Thus $l,r\in U$, a contradiction.
\end{proof}

\begin{lemma}[Inside-$U$ coset obstruction]\label{lem:coset}\status{PROVED}
Let $Q=\langle P_0,P\rangle\subseteq\Lambda^2(U)$, with $P$ rank two or redundant, and put $S=Q+\langle Z_1\rangle$. For $f,h\in V(S)$ and $g=fh$, put $W=V(S)+\langle g\rangle$. Then $z_2\notin W$, and no product of two elements of $W$ belongs to $z_2+W$.
\end{lemma}
\begin{proof}
Use the algebraic conclusions of [L3, B3a (Lemma 15)] and [L3, B3b (Lemma 16)] with their quadratic space replaced by $S$ and their first nonlinear product by $g$. This does not assume that $S$ is generated by three rank-two forms. The proofs of those two given algebraic results need only that the base space is $V(S)$ and $S\subseteq\Lambda^2(U)$. In B3a the coefficient vectors take values in $S$; the relevant kernel has dimension two because the external coefficient matrix has size $2\times4$ and rank two, independently of $\dim S$ and rank-two generation. B3b uses only constant external derivatives on $V(S)$ and B3a. The given arguments therefore apply unchanged to the present spaces. All 11 application spaces, including the ones not generated by rank-two forms, are independently checked below.

B3a first excludes $g\in z_2+V(S)$, hence $z_2\notin W$. A proposed final product either has both factors in $V(S)$, excluded by B3a even with a $g$ output correction, or normalizes to $k(g+F)$ with $k,F\in V(S)$. A constant change in $k$ absorbs a $g$ output correction and gives a product in $z_2+V(S)$, excluded by B3b.
\end{proof}

\begin{proposition}[One higher-degree signal is insufficient]\label{prop:onehigh}\status{PROVED}
Let $Q=\langle P_0,P\rangle$, with $P$ rank two or redundant, $g=fh$ with $f,h\in V(Q)$, and $S=Q+\langle Z_1\rangle$. Then $z_2\notin V(S)+\langle g\rangle$, and no product of two elements of this space places $z_2$ in it. The statement also allows $\deg g\le2$.
\end{proposition}
\begin{proof}
If $Q\subseteq\Lambda^2(U)$, apply Lemma~\ref{lem:coset}, since $f,h\in V(S)$.

Otherwise $P\notin\Lambda^2(U)$. The quartic part of $g$ is a scalar multiple of $P_0\wedge P$, which cannot equal $\Omega$ by Lemma~\ref{lem:rankext}. Thus $z_2$ is not already available.

If the proposed new product has both factors in $V(S)$, its quartic part is divisible by $P$: expand their quadratic parts in $P_0,P,Z_1$ and use $P_0\wedge Z_1=0$ and vanishing wedge squares. Adding a possible $g$ output correction retains this divisibility. Lemma~\ref{lem:rankext} excludes $\Omega$.

Otherwise normalize the factors and absorb the $g$ output correction to obtain
\[
k(g+F)=R\in z_2+V(S),\qquad k,F\in V(S).
\]
Since $kR=R$, Lemma~\ref{lem:lastfactor} gives $k=1$ or $k=1+P_0$. The first would put $z_2$ in the already excluded space $V(S)+\langle g\rangle$. In the second case reduction modulo $I$ gives
\[
g\equiv q_2+v+F\pmod I\qquad(v\in V(S)),
\]
because $1+P_0\equiv1$ and the carry belongs to $I$. Lemma~\ref{lem:ideal}, applied to $g=fh$, forces a form in $q_2+S$ to have rank at most two after adding an element of $Q$. All such forms have rank at least four by Lemma~\ref{lem:rankext}. This contradiction proves the proposition.
\end{proof}

\section[Proof of the main theorem]{Proof of Theorem~\ref{thm:main}}
The branch $\deg g_2\le2$ was excluded by Lemma~\ref{lem:replace} and [L3, Theorem 3]. Assume $\deg g_2\ge3$ and put $g=g_2$, $B=V(Q)$, $S=Q+\langle Z_1\rangle$.

Neither $z_2$ nor $z_1+z_2$ belongs to $V_4=B+\langle g,g_3\rangle$. Either placement, after freely adjoining $z_1$ in $V(S)$, would put $z_2$ in the span of $V(S)$, $g$ and one product with factors in $V(S)+\langle g\rangle$, contradicting Proposition~\ref{prop:onehigh}.

If $z_1\notin V_4$, the target cosets of $z_1$ and $z_2$ must be the same nonzero coset in the one-dimensional extension $V_5/V_4$. Their sum puts $z_1+z_2$ in $V_4$, impossible. Thus $z_1\in V_4$.

Since $g$ has degree at least three and $Z_1\notin Q$, we have $z_1\notin B+\langle g\rangle$. Here $Z_1\notin Q$ follows because $Z_1$ and $Z_1+P_0$ have rank four, whereas $P$ has rank two. Thus $g_3$ places $z_1$ and
\[
V_4=B+\langle g,z_1\rangle=V(S)+\langle g\rangle.
\]
The last gate now contradicts Proposition~\ref{prop:onehigh} again. This completes the human proof and closes the 2,064 difficult pairs by a common argument. \qed

\section{Computational checks}
All functions are stored as 64-bit ANF masks. Multiplication uses an involutory Boolean M\"obius transform, bitwise AND on truth tables, and the inverse transform. A separate monomial-OR multiplication cross-checks this algebra. No SAT solver or third-party package is used.

A factor plane is a two-dimensional subspace modulo constants. Its three nonzero vectors are ordered $i<j<i+j$. Changing constants or a plane basis changes the product only by an available function. Dependent factors add nothing. A space is always reduced as a space of actual functions; redundant signal coordinates are discarded.

For a quadratic space $B$ and a preceding signal $g$, every next gate either has both factors in $B$, or normalizes to $k(g+F)$ with $k,F\in B$. To test $k(g+F)\in T+B$ exactly, enumerate every $k$ satisfying the necessary identity $(1+k)(T+v)=0$ for some $v\in B$, then enumerate the complete image $kB$ modulo $B$. Constants in $k$ absorb the possible $g$ term in the output correction. This is an exact finite test, not a derivative approximation. The implementation also applies to bases containing cubic functions.

The full search checks both earlier targets $z_2,z_1+z_2$, both placements of $z_1$, and the final gate. If $z_1$ is placed by $g_3$, its resulting space is exactly $V(Q)+\langle g_2,z_1\rangle$, independent of the placing plane. If $z_1$ is placed by $g_2$, every subsequent nonlinear plane is tested. The last gate is allowed to be any product, a safe enlargement. Both complete normalized searches return no circuit.

\begin{center}\small
\begin{tabularx}{\textwidth}{@{}l X X@{}}\toprule
Item & Exact checked domain & Result\\\midrule
Lemma~\ref{lem:replace} & All 606 $Q$, all nonlinear planes; 83 symmetry representatives & 833,824 representative planes; 18,912 quadratic products; zero failures\\
Lemma~\ref{lem:sym} & All 32 maps, 1,024 compositions, 64 input points/map & 606 spaces, 83 orbits; target cosets preserved\\
Lemma~\ref{lem:degree} & All 887,857 affine and nonlinear planes, plus dependent factors & 237,568 quartic products; zero failures\\
Lemma~\ref{lem:four} & All $k,R$ in its domain; 68 representatives & 17,629,184 pairs; zero failures\\
Lemma~\ref{lem:lastfactor} & All $k,R$ in its domain; 67 representatives & 70,254,592 pairs; zero failures\\
Lemma~\ref{lem:two} & All 67 representatives and four $k$; arbitrary $g$ & 268 exact linear systems, covering all $2^{64}$ functions $g$; infeasible\\
Proposition~\ref{prop:reorder} & All applicable first-gate output spaces & Exact span identities verified\\
Lemma~\ref{lem:ideal} & All factor planes, affine and nonlinear; all quadratic representatives modulo $I$ & 887,857 planes; 270,385 quadratic quotient products; zero failures\\
Lemma~\ref{lem:rankext} & All 651 rank-two forms; all exterior multipliers & 2,604 rank tests; 616 full exterior image tests\\
Lemma~\ref{lem:coset} & All 11 needed inside-$U$ spaces, every first product and final product & 445,145 first planes; zero failures\\
Proposition~\ref{prop:onehigh} & All 606 initial spaces, all affine and nonlinear first products and final factors & 887,857 first planes, 740,610 output cosets; zero failures\\
Higher-degree exact search & 814,912 first-gate planes / 694,540 distinct cosets & 37,518 admissible cosets; no early or final target\\
Derivative plus $\Omega$ & All 37,518 admissible cosets, plane weights retained & 1,776 surviving cosets = 2,064 pairs; all fail exact test\\
Full $\lambda=2$ check & All 833,824 first nonlinear planes over representatives & No circuit; independent of [L3, Theorem 3]\\\bottomrule
\end{tabularx}\end{center}
Each executable check includes a planted mutation or a relaxed positive instance. The exact screens were independently compared with direct enumeration on all 1,740 cubic first-gate output cosets of the one-linear-gate case: 17,650,560 next-gate planes, with no discrepancy. Files, counts, timings and controls are in \texttt{checks/results.json} and per-orbit checkpoints. Arbitrary-function claims are checked on their full domain when possible by exact linear elimination; the two quadratic multiplier lemmas are checked on their complete stated application domains, not on arbitrary pairs of quadratic functions outside them.

\section{Status}
\begin{center}\begin{tabularx}{\textwidth}{@{}X l@{}}\toprule
Item & Status\\\midrule
All named lemmas and propositions, including ideal extraction & PROVED\\
Full human exclusion of higher-degree first gate & PROVED\\
Theorem~\ref{thm:main}, as a human-style proof & PROVED\\
Independent exhaustive exclusion, all normalized $\lambda=2$ circuits & PASS (computation)\\\bottomrule
\end{tabularx}\end{center}
PROVED means a written argument here plus a passing exhaustive check of its exact application domain. It does not certify human review.

\begin{thebibliography}{9}
\bibitem{L3} [L3] \emph{A human proof of the three-linear-gate case for the $3\times3$ low-half multiplier}, supplied \texttt{lambda3\_proof.tex}, unchanged copy in \texttt{references/}.
\bibitem{L4} [L4] \emph{A human proof of the four-linear-gate case for the $3\times3$ low-half multiplier}, supplied \texttt{lambda4\_proof (2).tex}, unchanged copy under the canonical filename in \texttt{references/}.
\end{thebibliography}
\end{document}
